\documentclass[11pt, letterpaper]{article}

% --- Idioma y codificación ---
\usepackage[spanish]{babel}
\usepackage[utf8]{inputenc}
\usepackage[T1]{fontenc}
\usepackage{lmodern}

% --- Matemáticas ---
\usepackage{amsmath, amssymb, amsthm}

% --- Formato ---
\usepackage[margin=2.5cm]{geometry}
\usepackage{parskip}
\usepackage{booktabs}
\usepackage{hyperref}
\hypersetup{colorlinks=true, linkcolor=blue!70!black, urlcolor=blue!70!black}

% --- Entornos de teoremas ---
\theoremstyle{definition}
\newtheorem{definition}{Definición}[section]
\theoremstyle{plain}
\newtheorem{theorem}{Teorema}[section]
\newtheorem{lemma}[theorem]{Lema}
\newtheorem{corollary}[theorem]{Corolario}
\theoremstyle{remark}
\newtheorem{remark}{Observación}[section]

\title{%
  T02 --- Demostración formal:\\[4pt]
  Kruskal produce un MST por la propiedad de corte
}
\author{Curso Linux--C--Rust --- B15 Estructuras de datos en C}
\date{}

\begin{document}
\maketitle
\tableofcontents

% ===================================================================
\section{Modelo}
% ===================================================================

\begin{definition}
Grafo $G = (V, E)$ no dirigido, conexo, con función de peso
$w: E \to \mathbb{R}$. $|V| = n$, $|E| = m$.
\end{definition}

\begin{definition}
\textbf{Árbol generador mínimo (MST).} Un subgrafo $T \subseteq E$
que es un árbol (conexo, acíclico) que contiene todos los vértices,
y cuyo peso total $w(T) = \sum_{e \in T} w(e)$ es mínimo.
\end{definition}

\textbf{Convención.} Asumimos que todos los pesos son distintos. (Si
hay empates, el resultado sigue siendo correcto pero el MST puede no
ser único; la prueba se adapta con una regla de desempate
consistente.)

\begin{definition}
\textbf{Corte.} Una partición $(S,\, V \setminus S)$ de $V$ con
$S \neq \emptyset$ y $S \neq V$. Una arista \textbf{cruza} el corte
si tiene un extremo en $S$ y otro en $V \setminus S$.
\end{definition}

% ===================================================================
\section{Lema fundamental: propiedad de corte}
% ===================================================================

\begin{lemma}[Propiedad de corte]\label{lem:cut}
Sea $(S,\, V \setminus S)$ un corte de $G$, y sea $e$ la arista de
\textbf{peso mínimo} que cruza el corte. Entonces $e$ pertenece a
todo MST de~$G$.
\end{lemma}

\begin{proof}
Por contradicción. Sea $T^*$ un MST que no contiene $e$. Sea
$e = (u, v)$ con $u \in S$ y $v \in V \setminus S$.

Como $T^*$ es un árbol generador, existe un camino único
$p: u \leadsto v$ en $T^*$. Este camino cruza el corte
$(S,\, V \setminus S)$ al menos una vez. Sea $e' = (x, y)$ una
arista de $p$ que cruza el corte (con $x \in S$,
$y \in V \setminus S$).

Consideremos $T' = T^* \setminus \{e'\} \cup \{e\}$:

\begin{enumerate}
  \item \textbf{$T'$ es conexo.} Eliminar $e'$ de $T^*$ desconecta
    $T^*$ en dos componentes: una conteniendo $x$ (dentro de $S$) y
    otra conteniendo $y$ (dentro de $V \setminus S$). Pero $u$ está
    en la componente de $x$ y $v$ en la de $y$, así que agregar
    $e = (u,v)$ reconecta las dos componentes.

  \item \textbf{$T'$ es acíclico.} $T'$ tiene $n - 1$ aristas
    (misma cantidad que $T^*$) y es conexo, por lo tanto es un
    árbol.

  \item \textbf{$w(T') < w(T^*)$.} Como $e$ es la arista de peso
    mínimo que cruza el corte y $e'$ también cruza el corte:
    $w(e) < w(e')$ (pesos distintos). Así:
    \[
      w(T') = w(T^*) - w(e') + w(e) < w(T^*)
    \]
\end{enumerate}

Esto contradice que $T^*$ es un MST.
\end{proof}

% ===================================================================
\section{Lema fundamental: propiedad de ciclo}
% ===================================================================

\begin{lemma}[Propiedad de ciclo]\label{lem:cycle}
Sea $C$ un ciclo en $G$, y sea $e$ la arista de \textbf{peso máximo}
en $C$. Entonces $e$ no pertenece a ningún MST.
\end{lemma}

\begin{proof}
Por contradicción. Supongamos que $e = (u, v)$ pertenece a un MST
$T^*$. Eliminar $e$ de $T^*$ desconecta el árbol en dos componentes
$S$ (conteniendo $u$) y $V \setminus S$ (conteniendo $v$).

Como $C$ es un ciclo que pasa por $e$, el camino
$C \setminus \{e\}: u \leadsto v$ conecta $S$ con
$V \setminus S$. Existe al menos una arista $e'$ en este camino que
cruza el corte $(S,\, V \setminus S)$, y $e' \neq e$.

Como $e$ es la arista de peso máximo en $C$ y $e' \in C$:
$w(e') < w(e)$. Entonces
$T' = T^* \setminus \{e\} \cup \{e'\}$ es un árbol generador con
$w(T') < w(T^*)$. Contradicción.
\end{proof}

% ===================================================================
\section{Demostración 1: Kruskal produce un MST}
% ===================================================================

\begin{theorem}\label{thm:kruskal}
El algoritmo de Kruskal retorna un MST.
\end{theorem}

\begin{proof}
Demostramos dos propiedades:

\textbf{(A) Kruskal produce un árbol generador.}

El algoritmo acepta una arista solo si sus extremos están en
componentes distintas (no crea ciclo). Así, el conjunto de aristas
aceptadas es siempre un bosque (acíclico). Como $G$ es conexo y
cada arista rechazada tiene ambos extremos en la misma componente,
al procesar todas las aristas, las $n$ componentes iniciales se
reducen a~1 (o el algoritmo se detiene al aceptar $n - 1$ aristas).
El resultado es un árbol generador. \checkmark

\textbf{(B) Cada arista aceptada pertenece a algún MST.}

Sea $e_k = (u, v)$ la $k$-ésima arista aceptada por Kruskal. Al
momento de aceptarla, $u$ y $v$ están en componentes distintas
$C_u$ y $C_v$ del bosque parcial. Consideremos el corte
$(C_u,\, V \setminus C_u)$.

La arista $e_k$ cruza este corte. Afirmamos que $e_k$ es la arista
de \textbf{peso mínimo} que cruza el corte:

Toda arista de peso menor que $w(e_k)$ ya fue procesada (las
aristas se procesan en orden ascendente de peso). Si alguna de esas
aristas cruzara el corte $(C_u,\, V \setminus C_u)$ actual,
tendría un extremo en $C_u$ y otro fuera. Pero al procesarla, fue
rechazada (pues ambos extremos estaban en la misma componente). Sin
embargo, las componentes solo se \textbf{fusionan}, nunca se
dividen: si ahora un extremo está en $C_u$ y el otro fuera, no
podrían haber estado en la misma componente antes.
\textbf{Contradicción.} Así que ninguna arista más liviana cruza el
corte actual.

Por el Lema~\ref{lem:cut} (propiedad de corte), $e_k$ pertenece a
todo MST. \checkmark

\textbf{(A) + (B):} Kruskal produce un árbol generador cuyas
$n - 1$ aristas pertenecen cada una a algún MST. Como los pesos son
distintos, el MST es único, y las $n - 1$ aristas forman ese MST.
\end{proof}

% ===================================================================
\section{Demostración 2: aristas rechazadas fuera del MST}
% ===================================================================

\begin{theorem}\label{thm:rejected}
Toda arista rechazada por Kruskal no pertenece a ningún MST.
\end{theorem}

\begin{proof}
Sea $e = (u, v)$ una arista rechazada. Al momento de procesarla,
$u$ y $v$ están en la misma componente $C$. Esto significa que
existe un camino $p: u \leadsto v$ formado por aristas ya aceptadas
(todas con peso $< w(e)$, pues fueron procesadas antes).

El camino $p$ junto con $e$ forma un ciclo. Como $e$ es la última
arista procesada en este ciclo (todas las del camino ya fueron
aceptadas), $e$ tiene el \textbf{peso máximo} del ciclo.

Por el Lema~\ref{lem:cycle} (propiedad de ciclo), $e$ no pertenece
a ningún MST.
\end{proof}

% ===================================================================
\section{Demostración 3: unicidad del MST con pesos distintos}
% ===================================================================

\begin{theorem}\label{thm:unique}
Si todos los pesos de aristas son distintos, el MST es único.
\end{theorem}

\begin{proof}
Por contradicción. Supongamos que existen dos MSTs distintos $T_1$
y $T_2$ con $w(T_1) = w(T_2)$. Sea $e$ la arista de
\textbf{menor peso} que está en uno pero no en el otro. Sin
pérdida de generalidad, $e \in T_1 \setminus T_2$.

Agregar $e$ a $T_2$ crea un ciclo $C$ en $T_2 \cup \{e\}$. Todas
las aristas de $C$ excepto $e$ pertenecen a $T_2$. Sea $e'$ una
arista de $C \setminus \{e\}$ que no está en $T_1$ (debe existir al
menos una, pues agregar $e$ a $T_1$ no crea ciclo pero a $T_2$ sí).
Entonces $w(e') > w(e)$ (porque $e$ era la de menor peso en la
diferencia simétrica $T_1 \triangle T_2$).

Pero $T_2' = T_2 \setminus \{e'\} \cup \{e\}$ sería un árbol
generador con:
\[
  w(T_2') = w(T_2) - w(e') + w(e) < w(T_2)
\]
Esto contradice que $T_2$ es un MST.
\end{proof}

% ===================================================================
\section{Demostración 4: complejidad
  \texorpdfstring{$O(m \log n)$}{O(m log n)}}
% ===================================================================

\begin{theorem}\label{thm:complexity}
El algoritmo de Kruskal ejecuta en $O(m \log n)$.
\end{theorem}

\begin{proof}
Desglosamos por fases:

\textbf{1) Ordenar aristas:} $O(m \log m)$. Como
$m \leq n(n-1)/2 < n^2$:
\[
  \log m \leq \log n^2 = 2\log n
\]
Así $O(m \log m) = O(m \log n)$.

\textbf{2) Operaciones Union-Find.} Se realizan a lo sumo $m$
operaciones \texttt{find} (dos por arista: una para cada extremo) y
a lo sumo $n - 1$ operaciones \texttt{union}. Con unión por rango y
compresión de caminos:
\[
  O\!\big((2m + n - 1) \cdot \alpha(n)\big)
  = O(m \cdot \alpha(n))
\]
donde $\alpha$ es la inversa de Ackermann ($\alpha(n) \leq 4$ para
todo $n$ práctico).

\textbf{3) Total:}
\[
  T = O(m \log n) + O(m \cdot \alpha(n)) = O(m \log n)
\]

El costo está dominado por el ordenamiento. Si las aristas llegan
pre-ordenadas, Kruskal corre en
$O(m \cdot \alpha(n)) \approx O(m)$.
\end{proof}

% ===================================================================
\section{Demostración 5: invariante del bosque parcial}
% ===================================================================

\begin{theorem}\label{thm:forest}
Sea $F_k$ el conjunto de aristas aceptadas tras procesar las
primeras $k$ aristas. Entonces $F_k$ es un subconjunto de algún MST.
\end{theorem}

\begin{proof}
Por inducción sobre $k$.

\textbf{Base ($k = 0$).} $F_0 = \emptyset \subseteq T^*$ para
cualquier MST $T^*$. \checkmark

\textbf{Paso inductivo.} Supongamos $F_{k-1} \subseteq T^*$ para
algún MST $T^*$. Sea $e_k$ la $k$-ésima arista procesada.

\emph{Caso 1: $e_k$ es rechazada.}
$F_k = F_{k-1} \subseteq T^*$. \checkmark

\emph{Caso 2: $e_k$ es aceptada.} Por el
Teorema~\ref{thm:kruskal}, $e_k$ es la arista de peso mínimo que
cruza algún corte.

Si $e_k \in T^*$: $F_k = F_{k-1} \cup \{e_k\} \subseteq T^*$.
\checkmark

Si $e_k \notin T^*$: agregar $e_k$ a $T^*$ crea un ciclo $C$.
Alguna arista $e'$ de $C$ cruza el mismo corte que $e_k$ y
$e' \in T^* \setminus F_{k-1}$. Definimos
$T^{**} = T^* \setminus \{e'\} \cup \{e_k\}$. Como
$w(e_k) \leq w(e')$ (por la propiedad de corte), $T^{**}$ es un
MST. Y $F_k = F_{k-1} \cup \{e_k\} \subseteq T^{**}$. \checkmark

En todos los casos, $F_k$ es subconjunto de algún MST.
\end{proof}

% ===================================================================
\section{Verificación numérica}
% ===================================================================

Grafo con 6 vértices y 9 aristas:

\begin{table}[h]
\centering
\begin{tabular}{ccccc}
\toprule
\textbf{Arista} & \textbf{Peso} & \textbf{¿Aceptada?}
  & \textbf{Corte respetado} & \textbf{Propiedad} \\
\midrule
$(4,5)$ & 1 & Sí & $(\{4\},\, \{0,1,2,3,5\})$
  & Mín.\ que cruza \checkmark \\
$(1,2)$ & 2 & Sí & $(\{1\},\, \{0,2,3,4,5\})$
  & Mín.\ que cruza \checkmark \\
$(0,3)$ & 3 & Sí & $(\{0\},\, \{1,2,3,4,5\})$
  & Mín.\ que cruza \checkmark \\
$(0,1)$ & 4 & Sí & $(\{0,3\},\, \{1,2,4,5\})$
  & Mín.\ que cruza \checkmark \\
$(2,4)$ & 5 & Sí & $(\{0,1,2,3\},\, \{4,5\})$
  & Mín.\ que cruza \checkmark \\
$(1,4)$ & 6 & No & ---
  & Ciclo: máx.\ peso \checkmark \\
\bottomrule
\end{tabular}
\end{table}

Peso del MST: $1 + 2 + 3 + 4 + 5 = 15$. \checkmark

Ninguna arista de menor peso cruza el corte indicado en cada paso,
y la arista rechazada es la de mayor peso en su ciclo. \checkmark

% ===================================================================
\section{Resumen de resultados}
% ===================================================================

\begin{table}[h]
\centering
\begin{tabular}{llp{5.5cm}}
\toprule
\textbf{Resultado} & \textbf{Técnica} & \textbf{Clave} \\
\midrule
Propiedad de corte
  & Contradicción: intercambiar aristas
  & $T' = T^* \setminus \{e'\} \cup \{e\}$, $w(T') < w(T^*)$ \\
Propiedad de ciclo
  & Contradicción: eliminar máxima
  & Arista de peso máximo en ciclo $\notin$ MST \\
Kruskal produce MST
  & Corte + orden ascendente
  & Cada arista aceptada es mín.\ que cruza un corte \\
Aristas rechazadas $\notin$ MST
  & Propiedad de ciclo
  & Máxima en su ciclo con aristas previas \\
$O(m \log n)$
  & Ordenamiento domina
  & Union-Find $O(m \cdot \alpha(n))$ es sublineal \\
Invariante del bosque
  & Inducción sobre aristas procesadas
  & $F_k \subseteq$ algún MST en todo momento \\
\bottomrule
\end{tabular}
\end{table}

\end{document}
